\section{Conclusions and Open Problems} %\todo{CONT HERE}


We established a hierarchy theorem for chain-of-thought reasoning, identifying the number of decoding steps as a fundamental computational resource. Our results raise an important open question of obtaining a hierarchy theorem that is as tight as classical time and space hierarchies in the Turing machine model. In particular, can we show that
\[
\CoT(r(n)) \subsetneq \CoT(r(n)\,\mathrm{polylog}(r(n)))?
\]

Our proof relies on simulations of transformers with Turing machines combined with padding. 
The known simulations incur quadratic blow-up, and it is unclear whether near-linear time simulations of $\CoT(r(n))$ are possible. In particular, can we show
\[
\CoT(r(n)) \subseteq \mathrm{TIME}(r(n)^{1+o(1)})?
\]
Moreover, due to the use of the padding technique, we only get Theorem \ref{thm:main theorem} for $r(n) \geq n$. The case $r(n) < n$ is an open question for saturated-attention transformers. We expect a positive answer.

Finally, the full computational power of chain-of-thought reasoning is not yet well understood and is a promising research theme. 

\begin{ack}
% We thank [names] for helpful discussions.

N. V. Vinodchandran and Thanh Le were supported in part by NSF awards 2342244, 2413848, and a UNL Grand Challenges Grant.
A. Pavan was supported in part by NSF awards 2342245 and 2413849.
\end{ack}

